\documentclass{article}
\usepackage[margin=1in]{geometry}

\begin{document}

\begin{table}[p]
\centering
\caption{Summary of atomic data by element\label{tab:atomdata-summary}}
\small
\renewcommand{\arraystretch}{1.12}
\begin{tabular*}{0.9\textwidth}{@{\extracolsep{\fill}}lcrr}
\hline
\textbf{Element} & \textbf{Ion Stages} & \textbf{Levels} & \textbf{Lines} \\
\hline
$\mathrm{H}_{1}$ & I & 30 & 419 \\
$\mathrm{Si}_{14}$ & I--VI & 925 & 16749 \\
$\mathrm{Fe}_{26}$ & I--IX & 2988 & 61247 \\
\textbf{Total} &  & \textbf{3943} & \textbf{78415} \\
\hline
\end{tabular*}

\medskip
\parbox{0.9\textwidth}{\footnotesize Ion stages with level and line data are given in spectroscopic notation.}
\end{table}

\end{document}
